\documentclass[12pt,a4paper]{article}
\usepackage{ugmath}
\usepackage{float}
\usepackage{placeins}
\usepackage[labelfont=bf,labelsep=space]{caption}
\newcommand{\paren}[1]{\left( #1 \right)}
\newcommand{\alumno}{Ricardo León Martínez}
\newcommand{\materia}{Fisica I}
\newcommand{\profesor}{Lucero Uscanga Aguilera}
\newcommand{\tarea}{Tarea 3}
\newcommand{\fecha}{6/3/2026}

\begin{document}
    \begin{center}
    {\large\textbf{UNIVERSIDAD DE GUANAJUATO}}\\[0.3cm]
    {\normalsize\textbf{DIVISIÓN DE CIENCIAS NATURALES Y EXACTAS}}\\
    {\normalsize\textbf{CAMPUS GUANAJUATO}}\\[1cm]

    {\Large\textbf{\tarea\ (\materia)}}\\[1cm]
    \end{center}

    \textbf{Nombre:} \alumno \hfill 
    \textbf{Fecha:} \fecha \hfill 
    \textbf{Calificación:} \rule{3cm}{0.4pt} \\[0.3cm]

    \begin{mdframed}[style=mdbluebox,frametitle={Ejercicio 1}]
        Un cuerpo se mueve a lo largo de una recta. Su aceleración está dada por $a=-2x$, donde $x$ está en pies y $a$ en pies sobre $s^{-2}$. Encontrar la relación entre velocidad y distancia,
        suponiendo que cuando $x=0$, $v=4$ pies $s^{-1}$.
    \end{mdframed}

    \begin{solucion}
        Recordemos que
        \begin{equation*}
            a=\frac{dv}{dt}.
        \end{equation*}
        Pero si la aceleración depende de la posición entonces
        \begin{equation*}
            a=\frac{dv}{dt}=\frac{dv}{dx}\frac{dx}{dt}=v\frac{dv}{dx}.
        \end{equation*}
        Por lo tanto,
        \begin{equation*}
            v\frac{dv}{dx}=-2x.
        \end{equation*}
        Reescribimos
        \begin{equation*}
            v\frac{dv}{dx}=-2x.
        \end{equation*}
        Integramos ambos lados
        \begin{equation*}
            \int vdv=\int-2xdx.
        \end{equation*}
        En consecuencia,
        \begin{equation*}
            \frac{v^{2}}{2}=-x^{2}+C.
        \end{equation*}
        Multiplicando por 2
        \begin{equation*}
            v^{2}=-2x^{2}+C_{1}.
        \end{equation*}
        Usamos la condición inicial
        \begin{equation*}
            v(0)=4.
        \end{equation*}
        Entonces $C_{1}=16$. Por lo tanto
        \begin{equation*}
            v^{2}=16-2x^{2}.
        \end{equation*}
        Así, dado que inicialmente $v\gt0$
        \begin{equation*}
            v(x)=\sqrt{16-2x^{2}}.
        \end{equation*}
    \end{solucion}

    \begin{mdframed}[style=mdbluebox,frametitle={Ejercicio 2}]
        Una piedra es lanzada verticalmente hacia arriba desde el techo de un edificio con
        una velocidad de $29.4ms^{-1}$. Otra piedra se deja caer $4s$ después que se lanza
        la primera. Demostrar que la primera pasará a la segunda exactamente $4s$ después
        de que se soltó la segunda.
    \end{mdframed}

    \begin{proof}
        La ecuación de posición es
        \begin{equation*}
            y_{1}(t)=v_{0}t-\frac{1}{2}gt^{2}.
        \end{equation*}
        Sustituyendo $v_{0}=29.4$
        \begin{equation*}
            y_{1}(t)=29.4t-4.9t^{2}
        \end{equation*}
        La segunda piedra se suelta en $t=4$. Entonces la posición de la primera en ese
        instante es
        \begin{equation*}
            y_{1}(4)=29.4(4)-4.9(4)^{2}=39.2m
        \end{equation*}
        Por lo tanto, cuando se suelta la segunda piedra, la primera está a $39.2m$ arriba
        del punto de lanzamiento. Ademas su velocidad instantanea es
        \begin{equation*}
            v_{1}(4)=29.4-9.8(4)=-9.8ms^{-1}
        \end{equation*}
        Es decir va de bajada. Definamos
        \begin{equation*}
            \tau:=t-4
        \end{equation*}
        como el tiempo medido desde que se suelta la segunda piedra. Para la primera piedra
        su posición inicial es
        \begin{equation*}
            y_{1}(0)=39.2
        \end{equation*}
        y su velocidad inicial es
        \begin{equation*}
            v_{1}(0)=-9.8
        \end{equation*}
        Por lo tanto
        \begin{equation*}
            y_{1}(\tau)=39.2-9.8\tau-4.9\tau^{2}.
        \end{equation*}
        Para la segunda piedra, se suelta desde el reposo en el origen
        \begin{equation*}
            y_{2}(\tau)=-4.9\tau^{2}.
        \end{equation*}
        Las piedras se cuentran cuando
        \begin{equation*}
            y_{1}(\tau)=y_{2}(\tau).
        \end{equation*}
        Entonces
        \begin{equation*}
            39.2-9.8\tau-4.9\tau^{2}=-4.9\tau^{2}.
        \end{equation*}
        En consecuencia
        \begin{equation*}
            39.2-9.8\tau=0.
        \end{equation*}
        Así,
        \begin{equation*}
            \tau=\frac{39.2}{9.8}=4.
        \end{equation*}
        Las piedras entonces, se encuentran cuando
        \begin{equation*}
            \tau=4s.
        \end{equation*}
    \end{proof}

    \begin{mdframed}[style=mdbluebox,frametitle={Ejercicio 3}]
        La posición de una particular como una función del tiempo esta dada por
        \begin{equation*}
            \textbf{r}=[(2t^{2}-7t)\textbf{i}-t^{2}\textbf{j}]m.
        \end{equation*}
        Encuentra
        \begin{enumerate}
            \item[(a)] La velocidad en $t=2s$
            \item[(b)] La acelaración en $t=5s$.
            \item[(c)] La velocidad promedio entre $t=1s$ y $t=3s$.  
        \end{enumerate}
    \end{mdframed}

    \begin{solucion}
        \begin{enumerate}
            \item[(a)] Velocidad en $t=2$.\\
            Por definición,
            \begin{equation*}
                v(t)=\frac{dr}{dt}.
            \end{equation*}
            Derivando componente a componente
            \begin{equation*}
                v(t)=\frac{d}{dt}[(2t^{2}-7t)\textbf{i}-t^{2}\textbf{j}].
            \end{equation*}
            Por lo tanto
            \begin{equation*}
                v(t)=(4t-7)\textbf{i}+2t\textbf{j}
            \end{equation*}
            Evaluando en $t=2$
            \begin{equation*}
                v(2)=\textbf{i}+4\textbf{j}ms^{-1}.
            \end{equation*}
            \item[(b)] Aceleración en $t=5s$.\\
            La aceleración es
            \begin{equation*}
                a(t)=\frac{d\textbf{v}}{dt}.
            \end{equation*}
            Derivando
            \begin{equation*}
                a(t)=4\textbf{i}+2\textbf{j}.
            \end{equation*}
            Por lo tanto
            \begin{equation*}
                a(5)=4\textbf{i}+2\textbf{j}ms^{-2}.
            \end{equation*}
            \item[(c)] Velocidad promedio entre $t=1s$ y $t=3s$.\\
            Por definición
            \begin{equation*}
                \textbf{v}_{prom}=\frac{\textbf{r}(3)-\textbf{r}(1)}{3-1}.
            \end{equation*}
            La posición en $t=1$ es
            \begin{equation*}
                \textbf{r}(1)=-5\textbf{i}+9\textbf{j}.
            \end{equation*}
            La posición en $t=3$ es
            \begin{equation*}
                \textbf{r}(3)=-3\textbf{i}+9\textbf{j}.
            \end{equation*}
            Por tanto el desplazamiento es
            \begin{align*}
                \textbf{r}(3)-\textbf{r}(1)&=(-3\textbf{i}+9\textbf{j})-(-5\textbf{i}+9\textbf{j})\\
                &=2\textbf{i}+8\textbf{j}.
            \end{align*}
            Así, la velocidad promedio es
            \begin{align*}
                \textbf{v}_{prom}&=\frac{2\textbf{i}+8\textbf{j}}{2}\\
                &=\textbf{i}+4\textbf{j}ms^{-1}.
            \end{align*}
        \end{enumerate}
    \end{solucion}

\end{document}